$\varphi$ ist wegen

\mavergleichskette
{\vergleichskette
{ - { \mathrm i}
}
{ \notin }{ {\mathbb H}
}
{  }{
}
{    }{
}
{   }{
}
}
{}{}{}
auf ${\mathbb H}$ definiert. Wir zeigen

\mavergleichskettedisp
{\vergleichskette
{ \betrag {  { \frac{  z-  { \mathrm i}  }{ z + { \mathrm i}  } }  }
}
{ <} { 1
}
{ } {
}
{ } {
}
{ } {
}
}
{}{}{,}
was zu

\mavergleichskettedisp
{\vergleichskette
{ \betrag {  z- { \mathrm i}  }
}
{ <} { \betrag { z + { \mathrm i}  }
}
{ } {
}
{ } {
}
{ } {
}
}
{}{}{}
äquivalent ist. Mit

\mavergleichskette
{\vergleichskette
{z
}
{ = }{ a+b { \mathrm i}
}
{  }{
}
{    }{
}
{   }{
}
}
{}{}{}
ist dies äquivalent zu

\mavergleichskettedisp
{\vergleichskette
{ a^2 +(b+1)^2
}
{ >} { a^2 +(b-1)^2
}
{ } {
}
{ } {
}
{ } {
}
}
{}{}{,}
was wiederum äquivalent zu

\mavergleichskettedisp
{\vergleichskette
{ b^2+2b+1
}
{ >} { b^2 -2b +1
}
{ } {
}
{ } {
}
{ } {
}
}
{}{}{}
ist, was wegen

\mavergleichskette
{\vergleichskette
{ b
}
{ > }{ 0
}
{  }{
}
{    }{
}
{   }{
}
}
{}{}{}
stimmt.

Für

\mavergleichskette
{\vergleichskette
{ \betrag {  w  }
}
{ < }{ 1
}
{  }{
}
{    }{
}
{   }{
}
}
{}{}{}
ist direkt

\mavergleichskette
{\vergleichskette
{ 1-w
}
{ \neq }{ 0
}
{  }{
}
{    }{
}
{   }{
}
}
{}{}{,}
daher ist $\psi$ definiert. Sei

\mavergleichskette
{\vergleichskette
{ w
}
{ = }{ c+d  { \mathrm i}
}
{  }{
}
{    }{
}
{   }{
}
}
{}{}{}
mit

\mavergleichskette
{\vergleichskette
{ c^2+d^2
}
{ < }{ 1
}
{  }{
}
{    }{
}
{   }{
}
}
{}{}{.}
Dann ist

\mavergleichskettealignhandlinks
{\vergleichskettealignhandlinks
{ { \frac{   (w+1) { \mathrm i}  }{  1-w  } }
}
{ =} { { \frac{   (c+1+d { \mathrm i} ) { \mathrm i}  }{  1-c -d { \mathrm i}  } }
}
{ =} { { \frac{   -d +(c+1) { \mathrm i}  }{  1-c -d { \mathrm i}  } }
}
{ =} { { \frac{   -d +(c+1) { \mathrm i}  }{  1-c -d { \mathrm i}  } } \cdot { \frac{   1-c+d { \mathrm i}  }{  1-c +d { \mathrm i}  } }
}
{ =} { { \frac{   -d (1-c) -d(c+1) + ( (1+c)(1-c)-d^2 ) { \mathrm i}  }{  (1-c)^2 +d^2  } }
}
}
{}
{}{.}
Da der Nenner eine positive reelle Zahl ist, ist der Imaginärteil dieser Zahl wegen

\mavergleichskettedisp
{\vergleichskette
{ 1-c^2-d^2
}
{ >0} {
}
{ } {
}
{ } {
}
{ } {
}
}
{}{}{}
positiv, also liegt das Bild von $\psi$ in ${\mathbb H}$.

Wir berechnen nun die Hintereinanderschaltungen. Es ist

\mavergleichskettealign
{\vergleichskettealign
{ \varphi( \psi(w))
}
{ =} { \varphi \left( \frac{   (w+1) { \mathrm i}  }{  1-w  } \right)
}
{ =} { { \frac{   { \frac{  (w+1) { \mathrm i}  }{  1-w  } } - { \mathrm i}  }{  { \frac{   (w+1) { \mathrm i}  }{  1-w  } } + { \mathrm i}  } }
}
{ =} { { \frac{   (w+1) { \mathrm i} - (1-w) { \mathrm i}  }{  (w+1) { \mathrm i} + (1-w) { \mathrm i}  } }
}
{ =} { { \frac{   2 w { \mathrm i}  }{  2 { \mathrm i}  } }
}
}
{
\vergleichskettefortsetzungalign
{ =} { w
}
{ } {}
{ } {}
{ } {}
}
{}{}
und

\mavergleichskettealign
{\vergleichskettealign
{ \psi( \varphi(z))
}
{ =} { \psi \left(  { \frac{   z- { \mathrm i}  }{  z + { \mathrm i}  } }  \right)
}
{ =} { { \frac{   \left(  { \frac{   z- { \mathrm i}  }{ z + { \mathrm i}  } } +1  \right)  { \mathrm i}  }{  1- { \frac{   z- { \mathrm i}  }{ z + { \mathrm i}  } }  } }
}
{ =} { { \frac{   ( z- { \mathrm i} + z+ { \mathrm i} ) { \mathrm i}  }{  z+ { \mathrm i} - (z- { \mathrm i} )  } }
}
{ =} { { \frac{   2 z { \mathrm i}  }{  2 { \mathrm i}  } }
}
}
{
\vergleichskettefortsetzungalign
{ =} { z
}
{ } {}
{ } {}
{ } {}
}
{}{.}

 [[Kategorie:Latexseite]]
